\documentclass{report}
\usepackage{graphicx} % Required for inserting images
\usepackage{amsmath}
\usepackage{amsthm}
\usepackage{amssymb}
\usepackage[a4paper, left=2cm, right=2cm, top=2cm, bottom=2cm]{geometry}
\usepackage[table]{xcolor} % Necessario per colorare righe/colonne
\usepackage{pifont}
\definecolor{SALZO}{rgb}{0.9,0.9,0.9}
\definecolor{lightblue}{rgb}{0.78, 0.94, 1.00}
\setcounter{secnumdepth}{1}
\usepackage{tikz} % Per disegnare il diagramma
\usepackage{hyperref}
\usepackage{bm}
\usepackage{cancel}
\usepackage[italian]{babel}
\usepackage{authblk}
\usepackage{pgfplots} % Pacchetto per i grafici
\pgfplotsset{compat=1.18} % Compatibilità con la versione
\usepackage{subcaption}
\usepackage{footnote}
\usepackage{listings}
\usepackage{tcolorbox}
\tcbuselibrary{listings}
\newtheoremstyle{nostyle}
  {0pt} % Spazio sopra
  {0pt} % Spazio sotto
  {\normalfont} % Corpo normale (non corsivo)
  {} % Nessuna indentazione
  {\bfseries} % Titolo in grassetto
  {.} % Punto dopo il titolo
  { } % Spazio dopo il titolo
  {} % Formato del titolo
% Per fare le matrici utilizzare \begin{pmatrix} a & b \\
%\usepackage{colortbl} % Per colorare celle e righe
\theoremstyle{nostyle}
\newtheorem{lemma}{Lemma}[section]
\newtheorem{theorem}{Teorema}[section]
\newtheorem{definition}{Definizione}[section]
\usepackage{enumerate}
\usepackage{nicematrix}
\usepackage{booktabs}
\usepackage{enumitem}
\usepackage{float}
\usepackage{xcolor}
\definecolor{verdeScuro}{RGB}{0, 100, 80}

\title{Sistemi di controllo - Formulario} % chktex 8
\author{Giovanni Esposito, Danilo Di Primio, Cristian Amendola} 
\affil{Università degli studi La Sapienza}
\affil{Ingegneria Informatica e Automatica - 2 anno} % chktex 8
\date{Febbraio - Maggio 2026} % chktex 8



\begin{document}
\maketitle
\tableofcontents

\section{Processo di un ingegnere automatico}
METTERE IMMAGINE
\subsection{Modellizzazione lineare:}
\[ \begin{cases}
    \dot{x}(t) = Ax(t) + Bu(t) + Pz(t) \\
    y(t) = Cx(t) + Du(t) + Qz(t)
\end{cases} \]

\subsection{Modellizzazione non lineare:}
\[ \begin{cases}
    x(t) = f\Big( x(t), u(t), z(t)   \Big) \\
    y(t) = g\Big( x(t), u(t), z(t)   \Big) 
\end{cases} \]

\subsection{Modellizzazione nel dominio di Laplace:}
\[  \begin{cases}
    x(s) = \left(sI-A \right)\left( Bu(s)+Pz(s)  \right) \\
    y(s) = P(s)u(s) + M(s)z(s)
\end{cases}   
\]
dove:
\begin{itemize}
    \item $P(s) = C(sI-A)^{-1} B +D$ 
    \item $M(s) = C(sI-A)^{-1} P+ Q$
\end{itemize}



\subsection{Modellizzazione nel dominio z:}

\section{Sistemi di controllo a controreazione}
METTERE IMMAGINE

\section{Funzioni di trasferimento}
Dato $\displaystyle F = PG$:
\subsection{Funzione di trasferimento del sistema complessivo}
\[ W(s) =  \frac{F}{1+F}  = \frac{N_F}{N_F + D_F}\]
\subsection{Funzione di trasferimento relativa al disturbo}
\[ W_z(s) = \frac{M}{1+ F} = \frac{N_M}{D_M} \ \frac{D_F}{N_F+D_F} \]
\subsection{Funzione di trasferimento relativa all'errore}
\[ W_e(s)  = \frac{1}{1+F} = \frac{D_F}{N_F+D_F}  \]




























\end{document}


